% ARQUITECTURA_AUTORAL_PREEXISTENTE: semicuadrado bidireccional de N42.
% FORMALIZACION_NUEVA: grupoide de pares y ponderacion por estabilizadores.
\subsubsection{Incidence marquée et normalisation bidirectionnelle}
\label{sec:ii09-normalizacion-bidireccional}

Les invariants du lecteur conservent la nature de l'objet dont ils
proviennent. L'intersection marquée $B_K=H_e\cap P_\pi$ a pour cardinal
$b=4$ ; l'entier $q=7$ est la coordonnée du pivot $p_\varphi$ dans la
numérotation orientée. L'ordre de détermination du système de Steiner
est $p=5$. Les supports hexadique et octadique ont pour cardinaux $h_6=6$
et $o_8=8$, et l'alphabet d'orientation tritique a pour cardinal $t=3$.
Un cardinal de support et une coordonnée de pivot sont utilisés par
leurs applications respectives, même lorsqu'ils sont tous deux publiés comme entiers.

La normalisation de l'autocroisement pentagonal admet la
réalisation explicite suivante. Soit $P$ le support de cinq positions de la clôture
et soit $X=P\times P$ son espace de paires ordonnées. L'inversion
bidirectionnelle sur cet espace est
\[
 \jmath:X\longrightarrow X,\qquad \jmath(x,y)=(y,x),
 \qquad\jmath^2=I.
\]
L'information de stabilisateur de chaque orbite est conservée. Pour une
orbite $[z]$, on définit son poids par
$1/|\operatorname{Stab}_{C_2}(z)|$. La somme de ces poids est la
cardinalité du groupoïde d'action $X/\!/C_2$.

\begin{proposition}[Mesure orbitale de l'autocroisement bidirectionnel]
\label{prop:ii09-semicuadrado-orbital}
Pour un support fini de cardinal $p$, l'inversion des paires ordonnées
détermine
\[
 |(P\times P)/\!/C_2|
 =\sum_{[z]\in(P\times P)/C_2}
      \frac1{|\operatorname{Stab}_{C_2}(z)|}
 =\frac{p^2}{2}.
\]
Pour $p=5$, l'autocroisement possède dix orbites libres et cinq orbites
fixes, et sa mesure orbitale est $10+5/2=25/2$.
\end{proposition}
\begin{proof}
Chaque paire diagonale $(x,x)$ forme une orbite fixe, dont le stabilisateur
est d'ordre deux ; il y en a $p$. Les $p(p-1)$ paires restantes se
regroupent deux par deux, avec un stabilisateur trivial. Leur contribution est
donc $p(p-1)/2$, et celle de la diagonale est $p/2$. La somme vaut
$p^2/2$. De manière équivalente, l'identité orbite--stabilisateur donne
$1/|\operatorname{Stab}(z)|=|[z]|/2$ ; la sommation sur toutes les orbites
produit $|P\times P|/2$.
\end{proof}

La mesure conserve explicitement la contribution des points fixes.
L'ensemble ordinaire des orbites a pour cardinal $p(p+1)/2$, qui vaut
quinze dans la clôture pentagonale. La mesure orbitale et ce cardinal sont
deux lectures déterminées de la même inversion. La normalisation
bidirectionnelle du lecteur utilise la première : les dix orbites libres
pèsent une unité et les cinq diagonales pèsent une demi-unité.

Cette réalisation précise le facteur $p^2/2$ employé dans le système
sectoriel suivant. L'affectation de cette contribution au secteur $23$,
son orientation négative et les deux autres règles d'incidence restent
spécifiées par le système sectoriel complet. La démonstration établit la
normalisation de l'autocroisement et maintient distincte la sélection des
fonctionnelles angulaires.
